\documentclass[11pt,a4paper]{article}
\usepackage[margin=2.5cm]{geometry}
\usepackage{amsmath,amssymb,amsthm}
\usepackage{booktabs}
\usepackage[T1]{fontenc}
\usepackage{lmodern}
\usepackage{microtype}
\usepackage[hidelinks]{hyperref}
\hypersetup{
  pdftitle={Bilinear Maps in Fewer AND Gates than Rank: A Polynomial Gap Between Multiplicative and Quadratic-Circuit Complexity over F2},
  pdfauthor={Yuki Obuchi},
  pdfsubject={Multiplicative complexity of bilinear maps over F2},
  pdfkeywords={multiplicative complexity, bilinear maps, bilinear rank, quadratic Boolean forms, XOR-AND circuits, quadratic circuit optimality, Boolean circuit complexity},
  pdfcreator={LaTeX},
}
\DeclareMathOperator{\MC}{MC}
\DeclareMathOperator{\qMC}{qMC}
\newtheorem{theorem}{Theorem}
\newtheorem{lemma}{Lemma}
\theoremstyle{remark}
\newtheorem*{remark}{Remark}
\newcommand{\F}{\mathbb F_2}
\title{Bilinear Maps in Fewer AND Gates than Rank:\\ A Polynomial Gap Between Multiplicative and\\ Quadratic-Circuit Complexity over $\mathbb F_2$}
\author{Yuki Obuchi}
\date{Research note, first published 2026-11-26\\ \small \url{https://github.com/ObuchiYuki/bilinear-mc-gap}}
\begin{document}
\maketitle
\begin{abstract}
For $n=4^N$, every bilinear map $\varphi:\F^n\times\F^n\to\F^n$ is computed by a circuit of AND-depth two with $2\cdot 12^N=2n^{\log_4 12}=O(n^{1.7925})$ AND gates, XOR gates being free. All gate inputs come from one fixed frame; only the XOR wiring depends on $\varphi$. The frame is a tensor power of a $(4,4,4)$ frame with twelve products and twelve second-layer gates, itself the product of a Karatsuba-type block and its dual; the saving comes from cancelling cubic cross terms using $x\cdot x=x$. By counting, almost every such map has bilinear rank at least $n^2/3$ and therefore needs at least $n^2/6$ AND gates in any quadratic (level-one) circuit. Consequently, for $N\ge 9$ almost all bilinear maps satisfy $\MC(\varphi)<\qMC(\varphi)$ with ratio at least $n^{0.2075}/12$, and for $N\ge 7$ almost all satisfy $\MC(\varphi)<R(\varphi)$. The gap is non-explicit. To the best of our knowledge it is new; it strengthens a single explicit separation $6<7$ from a companion note. This note has not been peer-reviewed.
\end{abstract}

\section{Definitions and results}

All arithmetic is over $\F$. Variables are Boolean, so $x_k^2=x_k$, and we compute Boolean functions, i.e.\ elements of $\F[x,y]/(x_k^2-x_k,\,y_s^2-y_s)$, not formal polynomials.

A bilinear map $\varphi:\F^{n_x}\times\F^{n_y}\to\F^{d}$ is given by its coefficient tensor:
\[
\varphi_a(x,y)=\sum_{k,s}\varphi[k,s,a]\,x_ky_s\qquad(a=1,\dots,d).
\]
Its rank $R(\varphi)$ is the least $r$ such that $\varphi$ is a sum of $r$ tensors $u\otimes v\otimes w$.

Circuits consist of fan-in-two AND gates, and XOR gates and constants, which are free.
\begin{itemize}
\item $\MC(\varphi)$ is the least number of AND gates in any circuit computing $\varphi_1,\dots,\varphi_d$.
\item $\qMC(\varphi)$ is the least number of AND gates in a \emph{quadratic circuit}. In such a circuit every AND gate multiplies two affine functions of the inputs, and each output is an affine combination of AND gates and inputs. These are the level-one circuits $L_1$ of Mirwald and Schnorr~\cite{MS92} and the quadratic circuits of Boyar and Find~\cite{BF18}.
\end{itemize}

\begin{theorem}[universal two-layer frame]\label{thm:frame}
Let $n=4^N$. Every bilinear map $\varphi:\F^n\times\F^n\to\F^n$ is computed by a circuit of AND-depth two with exactly $12^N+12^N=2n^{\log_4 12}$ AND gates ($\log_412=1.79248\ldots$). All linear forms entering the AND gates and all output assignments are independent of $\varphi$; only a $0/1$ wiring matrix depends on $\varphi$. For arbitrary $n$, $\MC(\varphi)\le 24\,n^{\log_4 12}$.
\end{theorem}

\begin{theorem}[polynomial gap]\label{thm:gap}
Let $n=4^N$. For all but a $2^{-n}$ fraction of the bilinear maps $\varphi:\F^n\times\F^n\to\F^n$, the following hold.
\begin{enumerate}
\item If $N\ge7$, then $\MC(\varphi)\le2\cdot12^N<n^2/3\le R(\varphi)$.
\item If $N\ge9$, then $\MC(\varphi)\le2\cdot12^N<n^2/6\le\qMC(\varphi)$.
\end{enumerate}
In particular, $\qMC(\varphi)/\MC(\varphi)\ge n^{2-\log_412}/12=n^{0.2075\ldots}/12$.
\end{theorem}

Theorem~\ref{thm:frame} is constructive. Theorem~\ref{thm:gap} combines it with a counting bound, so it gives no explicit map.

\section{Two-layer frames}

A \emph{frame} of shape $(n_x,n_y,d)$ with $m$ gates and $t$ products consists of
\begin{itemize}
\item linear forms $\xi_j$ in $x$ and $\eta_j$ in $y$ ($j\le t$), and
\item linear forms $\lambda_i$ in $x$ with output sets $c_i\subseteq\{1,\dots,d\}$ ($i\le m$).
\end{itemize}
A \emph{wiring} is a matrix $Q\in\F^{m\times t}$. The frame with wiring $Q$ is the circuit
\[
p_j=\xi_j(x)\,\eta_j(y),\qquad h_i=\lambda_i(x)\cdot\sum_jQ_{ij}\,p_j,\qquad \mathrm{out}_a=\sum_{i:\,a\in c_i}h_i,
\]
with exactly $t+m$ AND gates and AND-depth two. The frame is \emph{universal} if every bilinear map of its shape is computed by some wiring. Write $c_i[a]=1$ if $a\in c_i$, and put
\[
\Omega[k,l,s,a]=\sum_{i,j}Q_{ij}\,\lambda_i[k]\,c_i[a]\,\xi_j[l]\,\eta_j[s].
\]

\begin{lemma}\label{lem:crit}
The circuit computes $\varphi$ if and only if $\Omega[k,l,s,a]=\Omega[l,k,s,a]$ for all $k\neq l$, and $\Omega[k,k,s,a]=\varphi[k,s,a]$, for all $s,a$.
\end{lemma}
\begin{proof}
Formally, $\mathrm{out}_a=\sum_{k,l,s}\Omega[k,l,s,a]\,x_kx_ly_s$. As a Boolean function this equals
\[
\sum_{k,s}\Omega[k,k,s,a]\,x_ky_s+\sum_{k<l,\,s}\bigl(\Omega[k,l,s,a]+\Omega[l,k,s,a]\bigr)x_kx_ly_s .
\]
Distinct multilinear monomials are linearly independent as Boolean functions.
\end{proof}

Since $\Omega$ is linear in $Q$, a frame is universal as soon as every basis tensor $e_k\otimes e_s\otimes e_a$ has a wiring.

\section{Two small frames}

\begin{lemma}[type A]\label{lem:A}
The frame of shape $(2,2,2)$ below is universal, with $3$ gates and $4$ products:
\begin{itemize}
\item gates $(x_1;\{1\})$, $(x_2;\{2\})$, $(x_1+x_2;\{1,2\})$;
\item all products $x_ky_s$.
\end{itemize}
\end{lemma}
\begin{proof}
Write $\varphi_1=x_1\alpha_1(y)+x_2\beta_1(y)$ and $\varphi_2=x_1\alpha_2(y)+x_2\beta_2(y)$. Every bilinear $q_i=x_1a_i(y)+x_2b_i(y)$ is a wire sum. Then
\begin{align*}
x_1q_1+(x_1+x_2)q_3&=x_1(a_1+a_3)+x_2b_3+x_1x_2(b_1+a_3+b_3),\\
x_2q_2+(x_1+x_2)q_3&=x_1a_3+x_2(b_2+b_3)+x_1x_2(a_2+a_3+b_3).
\end{align*}
Take $a_3=\alpha_2$, $b_3=\beta_1$, $a_1=\alpha_1+\alpha_2$, $b_1=a_2=\alpha_2+\beta_1$, $b_2=\beta_1+\beta_2$.
\end{proof}

\begin{lemma}[type B]\label{lem:B}
The frame of shape $(2,2,2)$ below is universal, with $4$ gates and $3$ products:
\begin{itemize}
\item products $x_1y_1$, $x_2y_2$, $(x_1+x_2)(y_1+y_2)$;
\item gates $(x_u;\{a\})$ for $u,a\in\{1,2\}$.
\end{itemize}
\end{lemma}
\begin{proof}
The sum of the three products is $x_1y_2+x_2y_1$. The four monomials are
\[
x_1y_1=x_1\cdot x_1y_1,\qquad x_2y_2=x_2\cdot x_2y_2,
\]
\[
x_1y_2=x_1\cdot(x_1y_2+x_2y_1)+x_2\cdot x_1y_1,\qquad x_2y_1=x_2\cdot(x_1y_2+x_2y_1)+x_1\cdot x_2y_2 .
\]
For output $a$, sum these over the monomials of $\varphi_a$. The terms with multiplier $x_u$ combine into the gate $(x_u;\{a\})$.
\end{proof}

\section{Tensor products; proof of Theorem~\ref{thm:frame}}

The product $F_1\otimes F_2$ of frames of shapes $(n_1,s_1,d_1)$ and $(n_2,s_2,d_2)$ has shape $(n_1n_2,s_1s_2,d_1d_2)$, with
\begin{itemize}
\item gates $(\lambda_{i_1}\otimes\lambda_{i_2};\,c_{i_1}\times c_{i_2})$,
\item products $(\xi_{j_1}\otimes\xi_{j_2})(x)\cdot(\eta_{j_1}\otimes\eta_{j_2})(y)$.
\end{itemize}

\begin{lemma}\label{lem:tensor}
If $F_1$ and $F_2$ are universal, so is $F_1\otimes F_2$.
\end{lemma}
\begin{proof}
If $Q_1$ computes $\varphi_1$ and $Q_2$ computes $\varphi_2$, then the wiring $Q_1\otimes Q_2$ gives
\[
\Omega[(k_1,k_2),(l_1,l_2),(s_1,s_2),(a_1,a_2)]=\Omega_1[k_1,l_1,s_1,a_1]\,\Omega_2[k_2,l_2,s_2,a_2].
\]
This tensor is symmetric in its first two indices and has diagonal $\varphi_1\otimes\varphi_2$, so $Q_1\otimes Q_2$ computes $\varphi_1\otimes\varphi_2$ by Lemma~\ref{lem:crit}. Basis tensors of the product shape are products of basis tensors, so linearity finishes the proof.
\end{proof}

\begin{proof}[Proof of Theorem~\ref{thm:frame}]
Type~A $\otimes$ type~B is a universal $(4,4,4)$ frame with $12$ gates and $12$ products. Its $N$-th tensor power is universal of shape $(4^N,4^N,4^N)$ with $12^N$ gates and $12^N$ products. For general $n$ with $4^{N-1}<n\le4^N$, pad with zeros: $2\cdot12^N=24\cdot12^{N-1}<24\,n^{\log_412}$.
\end{proof}

\section{Lower bounds; proof of Theorem~\ref{thm:gap}}

\begin{lemma}\label{lem:q}
For every bilinear map $\varphi$, $\qMC(\varphi)\ge R(\varphi)/2$.
\end{lemma}
\begin{proof}
Let the AND gates be $g_r=A_rB_r$ with $A_r=\alpha_r(x)+\alpha'_r(y)+\epsilon_r$ and $B_r=\beta_r(x)+\beta'_r(y)+\epsilon'_r$ affine. A monomial $x_ky_s$ arises in $A_rB_r$ only as $x_k\cdot y_s$ or $y_s\cdot x_k$, and the reduction $x_k^2\to x_k$ does not touch it. So the $x$--$y$ part of $g_r$ is $\alpha_r\otimes\beta'_r+\beta_r\otimes\alpha'_r$.

Each output is an affine combination of gates and inputs. By uniqueness of the reduced polynomial, $\varphi$ is a sum of at most $2k$ rank-one tensors.
\end{proof}

\begin{lemma}\label{lem:count}
Let $n=4^N$. At most a $2^{-n}$ fraction of $n\times n\times n$ tensors over $\F$ have rank $<n^2/3$.
\end{lemma}
\begin{proof}
At most $2^{3nr}$ tensors have rank $\le r$. Since $n^2\equiv1\pmod3$, rank $<n^2/3$ means rank $\le r=(n^2-1)/3$, and then $3nr=n^3-n$.
\end{proof}

\begin{proof}[Proof of Theorem~\ref{thm:gap}]
By Theorem~\ref{thm:frame}, $\MC(\varphi)\le2\cdot12^N$. By Lemma~\ref{lem:count}, outside a $2^{-n}$ fraction we have $R(\varphi)\ge16^N/3$, and then $\qMC(\varphi)\ge16^N/6$ by Lemma~\ref{lem:q}.

The inequality $2\cdot12^N<16^N/3$ is equivalent to $(4/3)^N>6$, i.e.\ $N\ge7$. The inequality $2\cdot12^N<16^N/6$ is equivalent to $(4/3)^N>12$, i.e.\ $N\ge9$.
\end{proof}

\section{Remarks and scope}

\begin{remark}[counting lower bound for MC]
There are at most $2^{(2s+n)(s+2n+1)}$ circuits with $s$ AND gates. Hence most bilinear maps need $\MC=\Omega(n^{3/2})$, and the generic exponent lies between $3/2$ and $\log_412$. The best exponent of a universal frame is open.
\end{remark}

\begin{remark}[asymmetry is necessary]
Identify the $x$- and $y$-coordinates, and suppose every wire sum $q_i$ is a symmetric bilinear form. Put $S(u,v,w)=\sum_ic_i\lambda_i(u)q_i(v,w)$. If the output is bilinear, Lemma~\ref{lem:crit} makes $S$ symmetric in its first two arguments. Symmetry of the $q_i$ makes it symmetric in its last two arguments, so $S(x,x,y)=S(y,x,x)=\sum_ic_i\lambda_i(y)q_i(x,x)$. Every computed form is then a sum of at most $m$ rank-one forms, so such circuits never beat the rank.
\end{remark}

\begin{remark}[what is not claimed]
\begin{itemize}
\item Nothing is claimed about explicit maps.
\item Nothing is claimed about matrix multiplication: with $n=k^2$, Theorem~\ref{thm:frame} gives $O(k^{3.585})$.
\item Nothing is claimed about asymptotic rank.
\item Nothing is claimed about total circuit size. The wirings are dense, so the circuits have about $n^{3.585}$ XOR gates, and most maps need total size $\Omega(n^3/\log n)$ anyway.
\item For $N\le8$ the counting argument does not guarantee a gap, which does not mean that none exists.
\end{itemize}
\end{remark}

\begin{remark}[quadratic-circuit optimality]
Mirwald and Schnorr~\cite{MS92} asked whether $L=L_1$ for sets of quadratic forms; equality is known for one or two forms. An explicit set with $\MC=6<7=\qMC$ is given in~\cite{Ob26}. Theorem~\ref{thm:gap} shows that the ratio $L_1/L$ is unbounded, already for bilinear maps.
\end{remark}


\section*{Verification}
The repository contains \texttt{verify.py}, a standalone verifier that uses only the Python standard library (Python~3.10+). It performs the following checks:
\begin{itemize}
\item types~A and~B compute all 256 maps of their shape;
\item the $(4,4,4;12,12)$ base frame computes all 64 basis tensors on all 256 inputs, hence all $2^{64}$ maps by linearity, and its wiring is the unique solution of the linear system of Lemma~1 (rank 144);
\item the $(16,16,16;144,144)$ frame computes random dense targets, checked as exact reduced polynomials;
\item the integer thresholds $N\ge7$ and $N\ge9$ hold.
\end{itemize}
It runs automatically on every push.

\section*{Lean formalization}
Theorems~\ref{thm:frame} and~\ref{thm:gap} are proved in Lean~4 with Mathlib in the directory \texttt{lean/} of the repository (theorems \texttt{universal\_circuit}, \texttt{rank\_le\_two\_mul\_size}, \texttt{rank\_gap}, \texttt{gap}, \texttt{mc\_lt\_qmc}). There is no \texttt{sorry}, and the only axioms are \texttt{propext}, \texttt{Classical.choice} and \texttt{Quot.sound}. The statements use the definitions in \texttt{lean/BilinearMcGap/Defs.lean}, which the reader should compare with Section~1.

\section*{Acknowledgments}
The construction was found, and its proofs and verifier were drafted and checked, with the help of AI assistants (Claude by Anthropic; Gemini by Google for an independent cross-check). The author reviewed the final statements.

\begin{thebibliography}{9}
\bibitem{MS92} R.~Mirwald and C.-P.~Schnorr. The multiplicative complexity of quadratic Boolean forms. \emph{Theoretical Computer Science} 102(2):307--328, 1992. doi:10.1016/0304-3975(92)90235-8. Question in \S5, p.~325.
\bibitem{BF18} J.~Boyar and M.~G.~Find. Multiplicative complexity of vector valued Boolean functions. \emph{Theoretical Computer Science} 720:36--46, 2018. doi:10.1016/j.tcs.2018.02.023; arXiv:1407.6169.
\bibitem{Ob26} Y.~Obuchi. Six AND gates beat seven: a counterexample to quadratic-circuit optimality over $\mathbb F_2$. Research note, 2026. doi:10.5281/zenodo.24525440.
\end{thebibliography}


\vfill
\noindent\small This note is licensed under the Creative Commons Attribution 4.0 International License (CC BY 4.0). The accompanying code is MIT-licensed.
\end{document}
